\documentclass[11pt]{article}
\usepackage{fontspec}
\usepackage[french,english]{babel}
\usepackage{amsmath}
\usepackage{amssymb}
\usepackage{amsthm}
\usepackage{mathrsfs}
\usepackage{tikz-cd}
\usepackage{enumitem}

\newcommand{\printedpage}[1]{\marginpar{\footnotesize p.~#1}}
\newcommand{\msp}[1]{\marginpar{\footnotesize [#1]}}
\newcommand{\Cof}{\mathrm{Cof}}
\newcommand{\Fl}{\mathrm{Fl}}
\newcommand{\Mono}{\mathrm{Mono}}
\newcommand{\univ}{\mathrm{univ}}
\newcommand{\id}{\mathrm{id}}

% footnote marks in the text printed as superscript (n), in the footnote as superscript n
\makeatletter
\renewcommand\@makefnmark{\hbox{\textsuperscript{\normalfont(\@thefnmark)}}}
\renewcommand\@makefntext[1]{\parindent 1em\noindent\hbox{\textsuperscript{\normalfont\@thefnmark}}\,#1}
\makeatother

\begin{document}

\noindent\printedpage{14}% head: A. Grothendieck : Dérivateurs, Ch. XVI, Édition des Universités de Montpellier II et Paris VII
% foot: Version 18 mars 2019
\[
\begin{tikzcd}[ampersand replacement=\&, row sep=small, column sep=small]
X_0 \arrow[rr, "i_2"] \arrow[dr, "i_1"] \arrow[dd, "i_1"'] \& \& X_2 \arrow[dr, "j_1"] \arrow[dd, "j_1" near start] \& \\
\& X_1 \arrow[rr, "j_2" near start, crossing over] \& \& X_3 \arrow[dd, "\id"] \\
X_1 \arrow[rr, "j_2" near end] \arrow[dr, "\id"'] \& \& X_3 \arrow[dr, "\id"'] \& \\
\& X_1 \arrow[rr, "j_2"] \arrow[from=2-2, "\id"' near start, crossing over] \& \& X_3,
\end{tikzcd}
\tag*{(Q$'$)} % checked against PDF leaf 16: id labels the front vertical X_1 to X_1, not the horizontal X_1 to X_3.
\]
where $f_0 = i_1$, $f_1 = \id_{X_1}$, $f_1 = j_1$, $f_3 = \id_{X_3}$, hence % [sic] $f_1 = j_1$ -> $f_2 = j_1$, $f_1$ is already given as $\id_{X_1}$ and $f_2$ is missing
\[
j_1 \in W \Longrightarrow i_1 \in W, \ \ \textit{i.e.}\ \text{Loc (6 b)} \Longrightarrow \text{Loc (6$'$ b)},
\]
and if we have Loc (6 c), we see that $i_1 \in W \Longrightarrow j_1 \in W$, \textit{i.e.} Loc (6$'$ a).

\textbf{Comment 3.2.} Loc (6 a) appears to me as the ``vital minimum'' for establishing the existence of a $W$-homotopic theory of cofibers in Cat. For this, one must indeed have at hand a serviceable criterion ensuring that an arrow $i : A \longrightarrow B$ in Cat is in $\Cof_W$.

\msp{page 19}We find, when we have Loc (6 a) (in dual form, \textit{i.e.} for the \textit{closed ones,} it's all the same), that the \textit{open Dwyer immersions} are in $\Cof_W$, which suffices to have \textit{a good theory of $W$-homotopic cofibers, in which MV squares are $W$-cocartesian} (something we absolutely insist on). % [sic] a laquelle -> à laquelle, missing accent

As for Loc (6 b), if I have Loc (6$'$ b) + Loc (6 a) instead, I should manage to prove Loc (6). In XII Appendix (pages 265 ff.) I prove, it seems, that (with the present notation) Loc (6$'$) implies Loc (6), and even Loc (6 c) if $W$ is stable under direct factors (condition Loc (1 bis)). But the proof is incomplete, since I tacitly assume that the Dwyer immersions are in $\Cof_W$, something I had only established using axiom W (7), in fact using Loc (6 a) (which is part of W (7)).

\msp{page 20}But here the idea comes to me of the lemma of K. S. \textsc{Brown}, cited by \textsc{Baues}, which I am going to give again here:

\textbf{Lemma 3.2.1.} (K. S. \textsc{Brown}) \textit{Let $(\mathcal{M}, W)$ be a localizer} ($W$ \textit{mildly saturated}), \textit{equipped with a set $C \subset \Fl\,\mathcal{M}$ satisfying the conditions:}
\begin{enumerate}
\item[1)] \textit{$C$ stable under composition and under cobase change} (N.B. \textit{one has to assume the existence of the amalgamated sums}
$\begin{tikzcd}[ampersand replacement=\&, baseline=(\tikzcdmatrixname-1-1.base), row sep=small, column sep=small]
X_0 \arrow[rr, "i_1"] \arrow[dr, "i_2"'] \& \& X_1 \arrow[dr] \& \\
\& X_2 \arrow[rr] \& \& X_3
\end{tikzcd}$
\textit{only if $i_1 \in C$}), \textit{contains the isomorphisms.}
\item[2)] \textit{$\mathcal{M}$ contains an initial object, $\emptyset_{\mathcal{M}}$, and for every object $X$ of $\mathcal{M}$, $\emptyset_{\mathcal{M}} \longrightarrow X$ is in $C$.} (N.B. (1) \textit{and} (2) \textit{imply that $\mathcal{M}$ is stable under finite sums.})
\item[3)] \textit{Every arrow $f$ of $\mathcal{M}$ factors as $pi$, with $i \in C$, $p \in W$.}
\item[4)] \textit{$C \cap W$ is stable under cobase change,} i.e. \textit{if}
\[
\begin{tikzcd}[ampersand replacement=\&, row sep=small, column sep=small]
X_0 \arrow[rr, "i_1"] \arrow[dr, "i_2"'] \& \& X_1 \arrow[dr] \& \\
\& X_2 \arrow[rr, "j_1"'] \& \& X_3
\end{tikzcd}
\tag*{$(*)$}
\]
\printedpage{15}% head: A. Grothendieck : Dérivateurs, Ch. XVI, Édition des Universités de Montpellier II et Paris VII
% foot: Version 18 mars 2019
\textit{is cocartesian and $i_1 \in C \cap W$, then $j_1 \in C \cap W$.} (N.B. \textit{$j_1 \in C$ is already contained in} (1), \textit{it remains to require $j_1 \in W$.})
\end{enumerate}
\textit{Under these conditions, we have $C \subset \Cof_W$.}

One can also state the lemma as follows

\textbf{Corollary 3.2.2.} \textit{Suppose that $C$ satisfies} (1), (2), (3). \textit{Then}
\msp{page 21}\textit{the following conditions on $C$ are equivalent.}
\begin{enumerate}
\item[a)] \textit{Condition} (4) \textit{above,} i.e. \textit{stability of $C \cap W$ under cobase change.}
\item[a$'$)] \textit{$C \cap W \subset W^{\univ}$.}
\item[b)] \textit{$C \subset \Cof_W$.}
\item[b$'$)] \textit{For every cocartesian square}
\[
\begin{tikzcd}[ampersand replacement=\&, row sep=small, column sep=small]
X_0 \arrow[rr, "i"] \arrow[dr, "p"'] \arrow[to=2-4, phantom, "\textit{coc.}" description] \& \& X_1 \arrow[dr, "q"] \& \\
\& X'_0 \arrow[rr, "j"] \& \& X'_1
\end{tikzcd}
\]
\textit{with $i \in C$, $p \in W$, we have $q \in W$.}
\end{enumerate}
\textsc{Proof}. We have
\[
\begin{aligned}
&\text{b} \iff \text{b}'\\
&\text{a} \iff \text{a}'
\end{aligned}
\]
since $C$ is stable under cobase change. We have b $\Longrightarrow$ a$'$ since $\Cof_W \cap W = W^{\univ}$ thanks to hypothesis 3°). We have a $\Longrightarrow$ b$'$ by Brown's lemma.

We are going to apply this to the case $\mathcal{M} = \mathrm{Cat}$, $C =$ open Dwyer immersions. Conditions 1), 2) do not depend on $W$, they are satisfied, and 3) as well as soon as $W$ is a basic localizer. As for 4), it is a consequence of the old W (7 bis), \textit{i.e.} Loc (6$'$) (we
\msp{page 22}apparently need Loc (6$'$ a \textit{and} b).) We conclude $C \subset \Cof_W$ - it's wonderful! (But I don't know how to prove the lemma \dots)

This legitimizes the proof given in XII \textit{loc. cit.}, and we get the

\textbf{Proposition 3.3.} \textit{The condition} Loc (6$'$) (\textit{cocartesian square axiom, strong form} (a) \textit{and} (b)) is equivalent \textit{to} Loc (6 a, b) (\textit{axiom of homomorphisms of cocartesian squares}) \textit{and is} implied \textit{by} Loc (6 c). \textit{If $W$ satisfies} Loc (1 bis), i.e. \textit{is stable under direct factors, then} Loc (6$'$) is equivalent \textit{to} Loc (6 c) \textit{and to} Loc (6 a b c).

\textbf{Comment 3.4.} Thus, we get a theory of $W$-homotopic cofibers \textit{either} under hypothesis Loc (6$'$) (equivalent to the more complicated Loc (6)), or under the weaker (and more complicated) hypothesis Loc (6 a).

% left margin: a vertical line is printed alongside the following paragraph
In the first case, it is a theory of a very special nature, since the $W$-cocartesian squares are ``reversible'' (\textit{i.e.} we have not only $i_1 \in W \Longrightarrow j_1 \in W$, but also the reverse implication).

There remains then the question whether Loc (6$'$ a) suffices for a theory of $W$-homotopic cofibers, or what surely amounts to the same thing, whether it already implies Loc (6 a) (hence is equivalent to it). This hangs entirely on the \textit{question}:

\printedpage{16}% head: A. Grothendieck : Dérivateurs, Ch. XVI, Édition des Universités de Montpellier II et Paris VII
% foot: Version 18 mars 2019
\msp{page 23}% left margin: a question mark followed by a double vertical line alongside the first three lines of the following paragraph
\marginpar{\footnotesize ?\,$\Vert$}does the axiom Loc (6$'$ a) (cocartesian square axiom, direct form) imply that the set of open Dwyer immersions which are in $W$ ($W$-trivial Dwyer immersions) is stable under cobase change? In fact, we see that for this it would suffice to verify the axiom of the cube (or of homomorphisms of cocartesian squares), direct form Loc (6 a), in the case where $f_0$ and $f_2$ are identities, hence to prove that then $f_1 \in W \Longrightarrow f_3 \in W$ (which would directly give $\mathrm{Dw} \subset \Cof_W$). But I see no way of proving it \dots

Suppose now that condition Loc (6 a) is satisfied, so that we have a ``good theory of $W$-homotopic cofibers'', hence also a theory of $W$-cocartesian squares (\textit{cf.} XII, p. 258-262). There are two questions that arise, concerning the validity, with these weaker hypotheses, of the developments of XII:
\begin{enumerate}
\item[1°)] Given a commutative square
\[
\begin{tikzcd}[ampersand replacement=\&, row sep=small, column sep=small]
X_0 \arrow[rr] \arrow[dr] \& \& X_2 \arrow[dr] \& \\
\& X_1 \arrow[rr] \& \& X_3
\end{tikzcd}
\]
in Cat, is it true that this square is $W$-cocartesian if and only if the arrow
\[
\text{canonical}\ \displaystyle\int
\begin{tikzcd}[ampersand replacement=\&, row sep=small, column sep=small]
\& X_1 \\
X_0 \arrow[ur] \arrow[dr] \& \\
\& X_2
\end{tikzcd}
\quad \longrightarrow \ X_3 \text{ is in } W \text{ ? }\text{\footnotemark[15]}
\]
\footnotetext[15]{Better, in all cases, when we have a diagram
$\begin{tikzcd}[ampersand replacement=\&, baseline=(\tikzcdmatrixname-2-1.base), row sep=tiny, column sep=small]
\& X_1 \\
X_0 \arrow[ur] \arrow[dr] \& \\
\& X_2
\end{tikzcd}$, its integral $\displaystyle\int$ is a $W$-homotopic amalgamated sum of said diagram (assuming Loc (6 a) and no more, and even only Loc (6$'$ a)).}
\end{enumerate}
\msp{page 24}
\begin{enumerate}
\item[2°)] Is it still true that if $A$ is in Cat, and $u : F \lhook\joinrel\longrightarrow G$ a homomorphism in $A^{\wedge}$, then $u \in \Cof_{W_A}$, and if $u \in W$, $u \in W_A^{\univ}$, hence do we have
\[
\Mono_{A^{\wedge}} \subset \Cof_{W_A}, \qquad \Mono_{A^{\wedge}} \cap W_A \subset W_A^{\univ} \text{ ? }\text{\footnotemark[16]}
\]
\footnotetext[16]{N.B. Loc (6$'$ a) suffices in place of ( Loc (6 a)).}
\end{enumerate}
But 1°) is proved in XII, prop. 20, p. 262 - I have just checked that the proof is valid under the present hypotheses \footnote[17]{But it uses in addition Loc (3 bis). But with another simpler proof that I have just added (p 264), we can do without it (and we get a more complete view of the situation).}, provided that I check that corollary 3, p. 149 in XII is valid, \textit{i.e.} % [sic] (footnote 17) vison -> vision, missing letter
\begin{enumerate}
\item[1° bis)] If
$\begin{tikzcd}[ampersand replacement=\&, baseline=(\tikzcdmatrixname-1-1.base)]
X_0 \arrow[r] \arrow[d] \& X_1 \arrow[d] \\
X_2 \arrow[r] \& X
\end{tikzcd}$
is an MV square in Cat, then
$\displaystyle\int
\begin{tikzcd}[ampersand replacement=\&, row sep=small, column sep=small]
\& X_1 \\
X_0 \arrow[ur] \arrow[dr] \& \\
\& X_2
\end{tikzcd}
\quad \longrightarrow \quad X$ is in $W$ (hence, assuming 1°), the square is $W$-cocartesian).
\end{enumerate}

\printedpage{17}% head: A. Grothendieck : Dérivateurs, Ch. XVI, Édition des Universités de Montpellier II et Paris VII
% foot: Version 18 mars 2019
To check this, I am led to check that theorem 4 in loc. cit. (p. 146) is still valid.

I have just done so, and have annotated theorem 4 and its corollaries accordingly (some results (cor. 4 and cor. 5) using W (7 bis) in its full strength, \textit{i.e.} Loc (6$'$), \textit{i.e.} Loc (6 a \textit{and} b). But corollary 3 of theorem 4 is O.K.

I have also just checked the validity of 2°), by checking the proof of the corollary of theorem 5 (XII, p. 168) and annotating it accordingly.

\msp{page 25}3.5. Let us sum up: we have
\[
\text{Loc (6 c)} \ \Longrightarrow \ \text{Loc (6)} \ \iff \ \text{Loc (6$'$)}
\tag*{(3.5.1)}
\]
and
\[
\text{If } \ W = W^{\natural}\,, \ \text{ we have } \ \text{Loc (6 c)} \ \iff \ \text{Loc (6)} \ \iff \ \text{Loc (6$'$)}.
\tag*{(3.5.2)}
\]
On the other hand
\[
\text{Loc (6 a)} \ \Longrightarrow \ \text{Loc (6$'$ a)}\,,
\tag*{(3.5.3)}
\]
and I do not know whether the reverse implication
\[
\text{Loc (6$'$ a)} \ \overset{?}{\Longrightarrow} \ \text{Loc (6 a)}
\]
is valid \footnote[18]{This is the case in the presence of Loc (8, 9).}. To establish Loc (6 a), it suffices in any case to do so when $f_0$ and $f_2$ are isomorphisms.

The axiom Loc (6 a) seems to be the ``vital minimum'' for $W$ to give rise to a ``good theory of $W$-homotopic cofibers'', hence also of $W$-cocartesian squares. More precisely, I feel like proving this:

\textbf{Proposition 3.6.} \textit{The axiom} Loc (6 a) \textit{is equivalent to the conjunction of the following two properties.}

\printedpage{18}% head: A. Grothendieck : Dérivateurs, Ch. XVI, Édition des Universités de Montpellier II et Paris VII
% foot: Version 18 mars 2019
\begin{itemize}[leftmargin=7em, labelwidth=6em]
\item[\textbf{Loc (6 A)}] Cat has ``enough $W$-cofibrations'',
\msp{page 26}\textit{i.e.} every arrow $f$ of Cat factors as $f = pi$, with $i \in \Cof_W$ and $p \in W$. (This also means that Cat equipped with $(W, \Cof_W)$ is a category with cofibrations of K. S. \textsc{Brown} - \textsc{Anderson}, and all the results of \textsc{Brown-Quillen} theory are applicable. In particular, we have a good theory of $W$-cocartesian squares.)
\item[\textbf{Loc (6 B)}] If $X$ in Cat is the union of two opens $U$, $V$, with intersection $S$, then the square
\[
\begin{tikzcd}[ampersand replacement=\&, row sep=large, column sep=large]
S \arrow[r] \arrow[d] \& U \arrow[d] \\
V \arrow[r] \& X
\end{tikzcd}
\]
is $W$-cocartesian \footnote[19]{In fact, Loc (6 B) is a consequence of Loc (6 A), which therefore is \textit{equivalent} to Loc (6 a). (\textit{Cf.} sketch of proof in 4.4.15, pages 66-68.)}.
\end{itemize}
On the other hand, condition Loc (6 A) can be replaced by the condition (apparently weaker)
\begin{itemize}[leftmargin=7em, labelwidth=6em]
\item[\textbf{Loc (6 A$'$)}] Every Dwyer morphism is in $\Cof_W$.
\end{itemize}
\textbf{Corollary 3.7.} \textit{The condition} Loc (6), i.e. Loc (6 a, b) (\textit{or again the equivalent condition} Loc (6$'$)) \textit{is equivalent to the conjunction of conditions} Loc (6 A) (\textit{or} Loc (6 A$'$)) \textit{and} Loc (6 B) \textit{above}), \textit{and of the condition}
\begin{itemize}[leftmargin=7em, labelwidth=6em]
\item[\textbf{Loc (6 C)}] For every $W$-cocartesian square
\[
\begin{tikzcd}[ampersand replacement=\&, row sep=large]
X_0 \arrow[r, "i"] \arrow[d] \& X_1 \arrow[d] \\
X_2 \arrow[r, "j"] \& X
\end{tikzcd}
\]
in Cat, $j \in W \Longrightarrow i \in W$.
\end{itemize}
\msp{page 27}On the other hand, I suspect that Loc (6 c) is hardly ever satisfied unless $W = W^{\natural}$ (the most common case of course), in which case Loc (6 c) \textit{is equivalent} to Loc (6). There is probably no reason to include it in Loc (6), any more than $W = W^{\natural}$ in Loc (1).

\end{document}
